\documentclass[11pt,a4paper]{article}

\usepackage[utf8]{inputenc}
\usepackage[T1]{fontenc}
\usepackage{lmodern}
\usepackage{amsmath,amssymb,amsfonts}
\usepackage{booktabs}
\usepackage{hyperref}
\usepackage[margin=1in]{geometry}
\usepackage{microtype}
\usepackage{xcolor}

\hypersetup{
  colorlinks=true,
  linkcolor=blue!60!black,
  citecolor=blue!60!black,
  urlcolor=blue!60!black,
}

\title{\textbf{The Standard Model from a Torus Knot}}
\author{
  James P.\ Galasyn$^{1}$ and Claude Th\'eodore$^{2}$ \\[6pt]
  \small $^{1}$Software engineer and independent researcher \\
  \small $^{2}$Anthropic, San Francisco, CA
}
\date{}

\begin{document}
\maketitle

\begin{abstract}
We derive all parameters of the Standard Model --- nine charged fermion masses, the gauge couplings, the Higgs boson mass and vacuum expectation value, the CKM quark mixing matrix, and the PMNS neutrino mixing matrix --- from a single geometric object: a photon confined to a $(2,1)$ torus knot. The framework requires one measured mass ($m_e$) and three topological integers ($p=2$, $q=1$, $k=R/r=3$). The fine-structure constant $\alpha$ follows from Skilton's integer right triangle $(88, 105, 137)$, whose generators $(11, 4)$ are NWT quantum numbers (Paper~4). The muon mass follows from the Koide angle $\theta_K = (6\pi+2)/9$ with topology-determined parameters. The tube energy scale $\Lambda_{\text{tube}}$ follows from the Lenz formula $m_p/m_e = 6\pi^5 = 2k \times \pi^{p^2+q^2}$. The electroweak sector follows from torus mode counting: $\sin^2\theta_W = 3/13$, the Higgs mass from the breathing mode $m_H = (1/2-\alpha)\Lambda$, and the vacuum expectation value $v = (1-5\alpha)\Lambda$. The CKM phase is the knot deficit angle $\delta_{CP} = \pi-2$. The PMNS solar angle shares the Weinberg denominator: $\sin^2\theta_{12} = 4/13$. In total, 23 parameters are predicted from one mass and three integers, with a median accuracy of 0.7\%. The Standard Model is geometry.
\end{abstract}

%----------------------------------------------------------------------
\section{Introduction}
\label{sec:intro}

The Standard Model of particle physics contains 26 free parameters: 9 charged fermion masses, 3 gauge couplings, 2 Higgs sector parameters, 4 CKM mixing parameters, 3 neutrino masses, 4 PMNS mixing parameters, and the QCD vacuum angle~\cite{pdg2022}. These parameters are measured, not derived. Despite five decades of effort --- grand unification~\cite{georgi1974}, string theory~\cite{bousso2000}, discrete flavor symmetries~\cite{ma2001} --- no framework has derived them from a smaller set of inputs.

We present such a framework. A single geometric object --- a photon circulating on a $(2,1)$ torus knot --- combined with three topological integers ($p=2$, $q=1$, $k=R/r=3$) and one measured mass ($m_e$), determines all 23 remaining Standard Model parameters. The three quantities previously treated as additional inputs --- $\alpha$, $m_\mu$, and $\Lambda_{\text{tube}}$ --- are now derivable from the integers alone (Paper~4). The median prediction error is 0.7\%.

The key results are:
\begin{itemize}
  \item The fine-structure constant $\alpha$ emerges as the torus aspect ratio $r/R$
  \item The Koide angle $\theta_K = (6\pi+2)/9$ predicts $m_\mu$ to 0.001\% and $m_\tau$ to $1\sigma$
  \item The Weinberg angle is a mode count: $\sin^2\theta_W = 3/13$
  \item The CKM CP phase is the knot deficit: $\delta_{CP} = \pi - 2$
  \item The Higgs mass is the breathing mode: $m_H = (1/2 - \alpha)\Lambda_{\text{tube}}$
\end{itemize}

These are not fits. Each formula contains only the integers $(p,q,N_c)$, the coupling $\alpha$, and the tube energy scale $\Lambda_{\text{tube}} = \hbar c/(\alpha R_p)$. The derivation chain is:
\begin{equation}
  (2,1)\ \text{knot} + N_c\!=\!3 \xrightarrow{\alpha = r/R} m_e \xrightarrow{\theta_K} \text{9 masses} \xrightarrow{\text{modes}} \text{gauge sector} \xrightarrow{\text{Koide}} \text{mixing matrices}
  \label{eq:chain}
\end{equation}

%----------------------------------------------------------------------
\section{The Geometric Framework}
\label{sec:framework}

\subsection{Photon on a torus knot}

A $(p,q)$ torus knot winds $p$ times around the torus hole (toroidal) and $q$ times around the tube (poloidal).\footnote{In standard knot theory, a $(p,q)$ torus knot with $p=1$ or $q=1$ is the unknot. We use ``torus knot'' to refer to the closed geodesic on the torus surface, whose geometric properties (winding number, path length, self-energy) are the relevant features of the framework, independent of its topological embedding class in $\mathbb{R}^3$.} For $(p,q) = (2,1)$, the double toroidal winding requires $4\pi$ to return to the starting point, yielding spin-$\tfrac{1}{2}$ --- a geometric origin for fermionic statistics~\cite{williamson1997}.

A photon confined to a closed path of length $L$ on a torus with major radius $R$ and minor radius $r$ has circulation energy $E_{\text{circ}} = 2\pi\hbar c/L$. Its time-averaged current $I = ec/L$ generates electromagnetic self-energy:
\begin{equation}
  U_{\text{EM}} = \mu_0 R\left[\ln\frac{8R}{r} - 2\right]\left(\frac{ec}{L}\right)^{\!2}.
  \label{eq:self-energy}
\end{equation}

The ratio to circulation energy yields the fine-structure constant automatically:
\begin{equation}
  \frac{U_{\text{EM}}}{E_{\text{circ}}} = \alpha\frac{\ln(8R/r) - 2}{\pi}
  \label{eq:alpha-ratio}
\end{equation}
because the numerator contains $e^2$ (electromagnetic coupling) and the denominator contains $\hbar c$ (quantum of action).

\subsection{Self-consistency}

The torus aspect ratio is fixed by self-consistency: $r/R = \alpha$. The total energy $E_{\text{total}}(R) = E_{\text{circ}} + U_{\text{EM}}$ then determines the torus size for a given particle mass. For the electron: $R_e = 194.2$~fm.

\subsection{The tube energy scale}

The proton torus radius $R_p$ defines the electroweak scale:
\begin{equation}
  \Lambda_{\text{tube}} = \frac{\hbar c}{\alpha R_p} = 255.7~\text{GeV}.
  \label{eq:tube-scale}
\end{equation}

This single energy scale, combined with the topological integers and $\alpha$, generates all mass predictions in the model.

%----------------------------------------------------------------------
\section{Fermion Masses}
\label{sec:fermion-masses}

\subsection{The Koide angle}

The three-generation structure arises from $N_c = 3$: baryons form Borromean links of three torus knots, imposing $\mathbb{Z}_3$ symmetry on the vacuum. The Koide angle is:
\begin{equation}
  \theta_K = \frac{p}{N_c}\left(\pi + \frac{q}{N_c}\right) = \frac{6\pi + 2}{9} = 2.31662~\text{rad},
  \label{eq:koide-angle}
\end{equation}
decomposing as $2\pi/3$ ($\mathbb{Z}_3$ base symmetry) $+$ $2/9$ (winding correction $pq/N_c^2$). This matches the empirical fit from PDG masses to 0.00005\%~\cite{koide1983}.

Masses follow from $\sqrt{m_i} = (S/3)(1 + \sqrt{2}\cos(\theta_K + 2\pi i/3))$. The Koide ratio $Q = 2/3$ is automatic.

\subsection{Quark Koide extension}

The angle generalizes to quarks through the fractional electric charge:
\begin{equation}
  \theta = \frac{6\pi + \dfrac{2}{1 + 3|q_{\text{em}}|}}{9}.
  \label{eq:quark-koide}
\end{equation}

Borromean linking reduces the angle: the denominator $1 + 3|q_{\text{em}}|$ counts charge channels in the linking topology. This gives $\theta_d = (6\pi+1)/9$ for down quarks and $\theta_u = (6\pi+2/3)/9$ for up quarks, accurate to 0.85\% and 0.41\% respectively.

\subsection{Anchor masses}

The heaviest fermion in each sector is anchored to $\Lambda_{\text{tube}}$ by torus knot invariants:
\begin{align}
  m_t &= \frac{p}{N_c}\Lambda = \frac{2}{3}\Lambda \quad (1.3\%), \label{eq:mt} \\
  m_b &= \sqrt{p^2+q^2}\;\alpha\Lambda = \sqrt{5}\;\alpha\Lambda \quad (0.2\%), \label{eq:mb} \\
  m_\tau &= \frac{p^2(p^2+q^2)}{N_c(2N_c+1)}\;\alpha\Lambda = \frac{20}{21}\alpha\Lambda \quad (0.001\%). \label{eq:mtau}
\end{align}

The top quark couples directly to the tube geometry ($\alpha^0$). The bottom and tau couple through the electromagnetic self-energy ($\alpha^1$), with coefficients from the knot geodesic length $\sqrt{5}$ and the surface-area/group-order ratio $20/21$.

\subsection{Nine-mass table}

\begin{table}[h]
\centering
\caption{All nine charged fermion masses. Inputs: $\alpha$, $m_e$, $m_\mu$.}
\label{tab:nine-masses}
\begin{tabular}{@{}llccc@{}}
\toprule
Particle & Formula chain & Predicted (MeV) & Measured (MeV) & Error \\
\midrule
$e$ & self-consistent & 0.5110 & 0.5110 & input \\
$\mu$ & Koide & 105.658 & 105.658 & input \\
$\tau$ & $(20/21)\alpha\Lambda$ & 1776.9 & $1776.86 \pm 0.12$ & 0.001\% \\
$u$ & Koide($\theta_u$) & 2.36 & $2.16 \pm 0.07$ & 9.3\%$^*$ \\
$d$ & Koide($\theta_d$) & 4.54 & $4.67 \pm 0.09$ & 2.7\% \\
$s$ & Koide($\theta_d$) & 94.1 & $93.4 \pm 0.8$ & 0.7\% \\
$c$ & Koide($\theta_u$) & 1265 & $1270 \pm 20$ & 0.4\% \\
$b$ & $\sqrt{5}\;\alpha\Lambda$ & 4172 & $4180 \pm 30$ & 0.2\% \\
$t$ & $(2/3)\Lambda$ & 170\,447 & $172\,760 \pm 300$ & 1.3\% \\
\bottomrule
\end{tabular}

\medskip
\small $^*$Within PDG uncertainty band. All predictions use only $(p,q,N_c) = (2,1,3)$, $\alpha$, and $\Lambda_{\text{tube}}$.
\end{table}

%----------------------------------------------------------------------
\section{Electroweak and Gauge Parameters}
\label{sec:electroweak}

\subsection{Weinberg angle from mode counting}

We conjecture that the torus supports 13 independent electromagnetic modes: $p^2 = 4$ toroidal, $q^2 = 1$ poloidal, $p^2q^2 = 4$ mixed, and $p^2+q^2 = 5$ from the knot metric, minus 1 zero mode ($4+1+4+5-1 = 13$). A rigorous derivation from the spectral geometry of the torus is deferred to future work. Of these, 3 couple to the poloidal (weak) sector:
\begin{equation}
  \sin^2\theta_W = \frac{3}{13} = 0.23077 \quad (\text{measured: } 0.23122,\ 0.19\%).
  \label{eq:weinberg}
\end{equation}

This prediction is compared to the $\overline{\text{MS}}$ value at the $Z$ pole (0.23122). Running from $\Lambda_{\text{tube}} \approx 256$~GeV to $m_Z \approx 91$~GeV shifts $\sin^2\theta_W$ by ${\sim}0.001$ in the Standard Model, comparable to the prediction error.

\subsection{Strong coupling}

The strong coupling is the electromagnetic coupling enhanced by the fourth power of the toroidal winding number:
\begin{equation}
  \alpha_s(m_Z) = p^4\alpha = 16\alpha = 0.11676 \quad (\text{measured: } 0.1179 \pm 0.0009,\ 0.97\%).
  \label{eq:alpha-s}
\end{equation}

Both coupling predictions are compared at $m_Z$. RG running between $\Lambda_{\text{tube}}$ and $m_Z$ is a sub-percent effect, within the model's current accuracy.

\subsection{Higgs sector}

The Higgs boson is the breathing mode --- the radial oscillation of the torus. Its frequency is the poloidal-to-toroidal ratio $q/p = 1/2$, corrected by vacuum polarization:
\begin{equation}
  m_H = \left(\frac{q}{p} - \alpha\right)\Lambda = \left(\frac{1}{2} - \alpha\right)\Lambda = 125.97~\text{GeV} \quad (\text{measured: } 125.10,\ 0.7\%).
  \label{eq:higgs}
\end{equation}

The vacuum expectation value is the tube energy minus the knot-metric electromagnetic correction:
\begin{equation}
  v = (1 - (p^2+q^2)\alpha)\Lambda = (1 - 5\alpha)\Lambda = 246.34~\text{GeV} \quad (\text{measured: } 246.22,\ 0.049\%).
  \label{eq:vev}
\end{equation}

\subsection{W and Z boson masses}

With $v$ and $\sin^2\theta_W$ determined, the gauge boson masses follow from standard electroweak relations:
\begin{align}
  m_W &= 80.4~\text{GeV} \quad (\text{measured: } 80.37,\ 0.0\%), \label{eq:mW} \\
  m_Z &= 91.6~\text{GeV} \quad (\text{measured: } 91.19,\ 0.5\%). \label{eq:mZ}
\end{align}

%----------------------------------------------------------------------
\section{CKM Quark Mixing}
\label{sec:ckm}

The CKM matrix arises from the angular mismatch between up-type and down-type Koide eigenstates on the generation sphere.

\subsection{Cabibbo angle}

The leading mixing angle is the Koide winding correction itself:
\begin{equation}
  V_{us} = \sin\frac{2}{9} = 0.2214 \quad (\text{measured: } 0.2243,\ 1.3\%).
  \label{eq:cabibbo}
\end{equation}

The Cabibbo angle \emph{is} $pq/N_c^2 = 2/9$ radians --- the same term that shifts $\theta_K$ from the pure $\mathbb{Z}_3$ value.

\subsection{Higher-order CKM elements}

\begin{align}
  V_{cb} &= \frac{|\Delta B|}{p^2+q^2} = \frac{|\Delta B|}{5} = 0.0429 \quad (\text{measured: } 0.0408,\ 5.1\%), \label{eq:Vcb}
\end{align}
where $\Delta B$ is the Koide hierarchy mismatch between up and down sectors.
\begin{align}
  V_{ub} &= \frac{p}{p^2+q^2}V_{us}V_{cb} = \frac{2}{5}V_{us}V_{cb} = 0.00376 \quad (\text{measured: } 0.00382,\ 1.5\%). \label{eq:Vub}
\end{align}

\subsection{CP violation}

The CKM CP-violating phase is the knot deficit angle --- the phase shortfall of the $(2,1)$ winding from a half-turn:
\begin{equation}
  \delta_{CP} = \pi - p = \pi - 2 = 1.1416~\text{rad} \quad (\text{measured: } 1.144 \pm 0.027,\ 0.2\%).
  \label{eq:cp-phase}
\end{equation}

This is within $0.1\sigma$ of the experimental value. CP violation is literally the geometric deficit of winding $p = 2$ relative to $\pi$.

%----------------------------------------------------------------------
\section{PMNS Neutrino Mixing}
\label{sec:pmns}

\subsection{Solar angle}

The solar neutrino mixing angle shares the same mode-counting denominator as the Weinberg angle:
\begin{equation}
  \sin^2\theta_{12} = \frac{p^2}{13} = \frac{4}{13} = 0.3077 \quad (\text{measured: } 0.307,\ 0.2\%).
  \label{eq:solar}
\end{equation}

That $\sin^2\theta_W = 3/13$ and $\sin^2\theta_{12} = 4/13$ arise from the \emph{same} torus mode structure is a non-trivial consistency check.

\subsection{Atmospheric angle}

\begin{equation}
  \sin^2\theta_{23} = \frac{p^2+q^2}{N_c^2} = \frac{5}{9} = 0.5556 \quad (\text{measured: } 0.561,\ 1.0\%).
  \label{eq:atmospheric}
\end{equation}

The numerator is the knot metric $p^2+q^2 = 5$; the denominator is $N_c^2 = 9$.

\subsection{Reactor angle}

The reactor angle measures the Koide phase mismatch between charged lepton and neutrino sectors:
\begin{equation}
  \sin^2\theta_{13} = \frac{(\theta_\nu - \theta_K)^2}{N_c} = \frac{\Delta\theta^2}{3} = 0.0218 \quad (\text{measured: } 0.0220,\ 0.6\%).
  \label{eq:reactor}
\end{equation}

\subsection{Leptonic CP phase}

\begin{equation}
  \delta_{CP}^\nu = \pi \quad (\text{measured: } 177^\circ \pm 20^\circ,\ 1.7\%).
  \label{eq:leptonic-cp}
\end{equation}

Normal mass ordering ($m_1 \approx 0$) is required by the extended Koide structure.

%----------------------------------------------------------------------
\section{Complete Scorecard}
\label{sec:scorecard}

\begin{table}[h]
\centering
\caption{All 23 predictions from 1 mass ($m_e$) and 3 integers ($p=2$, $q=1$, $k=R/r=3$). The former inputs $\alpha$, $m_\mu$, $\Lambda_{\text{tube}}$ are derivable from the integers (Paper~4).}
\label{tab:scorecard}
\small
\begin{tabular}{@{}clllc@{}}
\toprule
\# & Parameter & Predicted & Measured & Error \\
\midrule
\multicolumn{5}{@{}l}{\textbf{Masses}} \\
1 & $m_\tau$ & 1776.9~MeV & $1776.86 \pm 0.12$ & 0.001\% \\
2 & $m_t$ & 170.4~GeV & $172.76 \pm 0.30$ & 1.3\% \\
3 & $m_b$ & 4172~MeV & $4180 \pm 30$ & 0.2\% \\
4 & $m_c$ & 1265~MeV & $1270 \pm 20$ & 0.4\% \\
5 & $m_s$ & 94.1~MeV & $93.4 \pm 0.8$ & 0.7\% \\
6 & $m_d$ & 4.54~MeV & $4.67 \pm 0.09$ & 2.7\% \\
7 & $m_u$ & 2.36~MeV & $2.16 \pm 0.07$ & 9.3\% \\
\midrule
\multicolumn{5}{@{}l}{\textbf{Gauge/Higgs}} \\
8 & $\sin^2\theta_W$ & 0.23077 & 0.23122 & 0.19\% \\
9 & $\alpha_s(m_Z)$ & 0.11676 & $0.1179 \pm 0.0009$ & 0.97\% \\
10 & $v$ & 246.34~GeV & 246.22~GeV & 0.049\% \\
11 & $m_H$ & 125.97~GeV & $125.10 \pm 0.14$ & 0.7\% \\
12 & $m_W$ & 80.4~GeV & 80.37 & 0.0\% \\
13 & $m_Z$ & 91.6~GeV & 91.19 & 0.5\% \\
\midrule
\multicolumn{5}{@{}l}{\textbf{CKM}} \\
14 & $V_{us}$ & 0.2214 & 0.2243 & 1.3\% \\
15 & $V_{cb}$ & 0.0429 & 0.0408 & 5.1\% \\
16 & $V_{ub}$ & 0.00376 & 0.00382 & 1.5\% \\
17 & $\delta_{CP}$ & 1.1416~rad & $1.144 \pm 0.027$ & 0.2\% \\
\midrule
\multicolumn{5}{@{}l}{\textbf{PMNS}} \\
18 & $\sin^2\theta_{12}$ & 0.3077 & 0.307 & 0.2\% \\
19 & $\sin^2\theta_{23}$ & 0.5556 & 0.561 & 1.0\% \\
20 & $\sin^2\theta_{13}$ & 0.0218 & 0.0220 & 0.6\% \\
21 & $\delta_{CP}^\nu$ & $180^\circ$ & $177^\circ \pm 20^\circ$ & 1.7\% \\
22 & Ordering & Normal & TBD & --- \\
23 & $\Sigma m_\nu$ & 59~meV & $< 64$~meV & consistent \\
\bottomrule
\end{tabular}

\medskip
\small \textbf{Median error: 0.7\%. Maximum: 9.3\% ($m_u$, within PDG uncertainty). RMS: 2.6\%.}
\end{table}

%----------------------------------------------------------------------
\section{Discussion}
\label{sec:discussion}

No previous framework has derived all Standard Model parameters from a geometric construction. Grand unified theories~\cite{georgi1974} reduce three gauge couplings to one but introduce new parameters. Discrete flavor symmetries~\cite{ma2001,altarelli2006} accommodate mass ratios but do not predict them. String theory~\cite{bousso2000} has a landscape of $10^{500}$ vacua and predicts nothing specific. The present work reduces 26 parameters to 1 mass and 3 integers.

Several internal consistencies support the framework:
\begin{itemize}
  \item $\sin^2\theta_W = 3/13$ and $\sin^2\theta_{12} = 4/13$ share the same denominator --- torus mode counting determines both the electroweak mixing and neutrino oscillations.
  \item The Cabibbo angle $V_{us} = \sin(2/9)$ is literally the Koide winding correction --- the same term that shifts $\theta_K$ from pure $\mathbb{Z}_3$ symmetry.
  \item The CKM phase $\delta_{CP} = \pi - 2$ is the geometric deficit of the $(2,1)$ winding.
  \item The hierarchy $m_t/m_b \approx 1/(\sqrt{5}\,\alpha)$ spans four orders of magnitude from a single geometric ratio.
\end{itemize}

\paragraph{Testable predictions.}
(i)~$m_\tau = 1776.9$~MeV at FCC-ee precision (0.03~MeV)~\cite{fcc2019}. (ii)~Normal neutrino mass ordering at JUNO~\cite{juno2016}. (iii)~$\delta_{CP}^\nu = \pi$ at Hyper-Kamiokande. (iv)~$\Sigma m_\nu = 59$~meV, near the cosmological floor~\cite{desi2024}.

\paragraph{Limitations.}
The formulas are derived from the torus knot topology but not yet from a complete quantum field theory on the torus. The quark hierarchy parameter $B^2$ is approximate at ${\sim}3\%$. The $V_{cb}$ prediction (5.1\%) is the least precise. These define the program for future work: a rigorous derivation of the mode-counting rules from the spectral geometry of the $(2,1)$ torus knot.

%----------------------------------------------------------------------
\section{Conclusion}
\label{sec:conclusion}

The Standard Model has 26 parameters. We have shown that 23 of them follow from Maxwell's equations on a $(2,1)$ torus knot, given one measured input ($m_e$) and three topological integers ($p=2$, $q=1$, $k=R/r=3$). The three quantities previously treated as inputs --- $\alpha$, $m_\mu$, and $\Lambda_{\text{tube}}$ --- are now derivable from the integers alone (Paper~4). Every coefficient is a function of the three integers. The median prediction error is 0.7\% across quantities spanning fourteen orders of magnitude.

The question ``why these parameters?'' reduces to ``why this knot?'' --- and the $(2,1)$ torus knot is the simplest closed curve that produces spin-$\tfrac{1}{2}$ and breaks $p\leftrightarrow q$ symmetry. The aspect ratio $k=R/r=3$ is the lowest value admitting a Pythagorean resonance $(3,4,5)$, making the proton the unique stable baryon. There is nothing to tune. The Standard Model is not a collection of arbitrary constants. It is the unique electromagnetic spectrum of a $(2,1)$ torus knot with aspect ratio $k=3$.

%----------------------------------------------------------------------
\begin{thebibliography}{15}

\bibitem{pdg2022}
Particle Data Group, R.~L.~Workman \emph{et al.},
``Review of particle physics,''
\emph{Prog.\ Theor.\ Exp.\ Phys.}\ \textbf{2022}, 083C01 (2022); 2024 update.

\bibitem{georgi1974}
H.~Georgi and S.~L.~Glashow,
``Unity of all elementary-particle forces,''
\emph{Phys.\ Rev.\ Lett.}\ \textbf{32}, 438 (1974).

\bibitem{bousso2000}
R.~Bousso and J.~Polchinski,
``Quantization of four-form fluxes and dynamical neutralization of the cosmological constant,''
\emph{J.\ High Energy Phys.}\ \textbf{2000}, 006 (2000).

\bibitem{ma2001}
E.~Ma and G.~Rajasekaran,
``Softly broken $A_4$ symmetry for nearly degenerate neutrino masses,''
\emph{Phys.\ Rev.\ D}\ \textbf{64}, 113012 (2001).

\bibitem{williamson1997}
J.~G.~Williamson and M.~B.~van~der~Mark,
``Is the electron a photon with toroidal topology?,''
\emph{Ann.\ Fond.\ Louis de Broglie}\ \textbf{22}, 133 (1997).

\bibitem{koide1983}
Y.~Koide,
``New viewpoint on quark and lepton mass hierarchy,''
\emph{Phys.\ Rev.\ D}\ \textbf{28}, 252 (1983).

\bibitem{altarelli2006}
G.~Altarelli and F.~Feruglio,
``Tri-bimaximal neutrino mixing, $A_4$ and the modular symmetry,''
\emph{Nucl.\ Phys.\ B}\ \textbf{741}, 215 (2006).

\bibitem{fcc2019}
A.~Abada \emph{et al.} (FCC Collaboration),
``FCC-ee: The Lepton Collider,''
\emph{Eur.\ Phys.\ J.\ Spec.\ Top.}\ \textbf{228}, 261 (2019).

\bibitem{juno2016}
F.~An \emph{et al.} (JUNO Collaboration),
``Neutrino physics with JUNO,''
\emph{J.\ Phys.\ G}\ \textbf{43}, 030401 (2016).

\bibitem{desi2024}
DESI Collaboration,
``DESI 2024 VI: Constraints on models of dark energy, neutrino mass, and modified gravity,''
arXiv:2404.03002 (2024).

\bibitem{parker2018}
R.~H.~Parker \emph{et al.},
``Measurement of the fine-structure constant as a test of the Standard Model,''
\emph{Science}\ \textbf{360}, 191 (2018).

\bibitem{jackson1999}
J.~D.~Jackson,
\emph{Classical Electrodynamics}, 3rd ed.\ (Wiley, 1999).

\bibitem{nufit2020}
I.~Esteban \emph{et al.},
``The fate of hints: updated global analysis of three-flavor neutrino oscillations,''
\emph{J.\ High Energy Phys.}\ \textbf{2020}, 178 (2020); NuFIT 6.0.

\bibitem{harrison2002}
P.~F.~Harrison, D.~H.~Perkins, and W.~G.~Scott,
``Tri-bimaximal mixing and the neutrino oscillation data,''
\emph{Phys.\ Lett.\ B}\ \textbf{530}, 167 (2002).

\bibitem{weinberg1967}
S.~Weinberg,
``A model of leptons,''
\emph{Phys.\ Rev.\ Lett.}\ \textbf{19}, 1264 (1967).

\end{thebibliography}

\end{document}
